\documentclass[12pt,a4paper]{article}
\usepackage[utf8]{inputenc}
\usepackage[T1]{fontenc}
\usepackage[english]{babel}
\usepackage{geometry}
\usepackage{hyperref}
\usepackage{booktabs}
\usepackage{longtable}
\usepackage{array}
\usepackage{listings}
\usepackage{xcolor}
\usepackage{parskip}
\usepackage{fancyhdr}
\usepackage{amsmath}
\usepackage{tcolorbox}
\usepackage{enumitem}
\usepackage{ragged2e}
\tcbuselibrary{skins,breakable}

\geometry{margin=2.5cm}
\hypersetup{colorlinks=true, linkcolor=blue!60!black, urlcolor=blue!60!black}

\definecolor{codebg}{RGB}{245,245,245}
\definecolor{noteblue}{RGB}{220,235,250}
\definecolor{warnred}{RGB}{255,235,230}
\definecolor{importantbg}{RGB}{230,250,235}
\definecolor{commentgreen}{RGB}{0,128,0}
\definecolor{keywordblue}{RGB}{0,0,180}
\definecolor{stringbrown}{RGB}{160,60,0}

\lstset{
  backgroundcolor=\color{codebg},
  basicstyle=\ttfamily\small,
  keywordstyle=\color{keywordblue}\bfseries,
  commentstyle=\color{commentgreen}\itshape,
  stringstyle=\color{stringbrown},
  breaklines=true, breakatwhitespace=true,
  frame=single, rulecolor=\color{gray!40},
  xleftmargin=6pt, xrightmargin=6pt,
  showstringspaces=false, language=Python
}

\newcolumntype{P}[1]{>{\RaggedRight\arraybackslash}p{#1}}

\pagestyle{fancy}
\fancyhf{}
\rhead{\texttt{reasoner.py} Documentation}
\lhead{\leftmark}
\cfoot{\thepage}

\title{\textbf{Rientra@ Semantic Microservice}\\[0.4em]
       \large Logic Layer Documentation\\[0.4em]
       \normalsize \texttt{Graphic/rientra-app/python-service/reasoner.py}}
\author{STIIMA-CNR}
\date{Version 3.1 --- May 2026}

\begin{document}
\begin{titlepage}\centering\maketitle\end{titlepage}
\tableofcontents
\newpage

\clearpage
\section{Overview}

\texttt{reasoner.py} is the \textbf{pure logic layer} of the Rientra@ semantic microservice. It wraps \texttt{owlready2} and the Pellet DL reasoner, exposing a set of public query and mutation functions that the FastAPI layer (\texttt{main.py}) calls from HTTP endpoints. It produces \textbf{no terminal output, no Rich tables, and no matplotlib figures}.

\subsection{Key Design Properties}

\begin{itemize}[noitemsep]
  \item \textbf{Separation of concerns}: all OWL/SPARQL logic is isolated here; \texttt{main.py} only handles HTTP routing.
  \item \textbf{Thread safety}: a \texttt{threading.Lock} serialises \texttt{isSelected} mutations; a second lock serialises health-condition mutations.
  \item \textbf{Two-tier serving}: immediately after startup a local JSON snapshot is served (fast); in the background Pellet runs and once complete the live ontology is used instead.
  \item \textbf{Snapshot caching}: after every Pellet run the full inferred read-model is serialised to \texttt{reasoning\_cache/inference\_\{fingerprint\}.json} so the next startup is instant.
  \item \textbf{Graceful degradation}: Pellet runs in a background thread; the HTTP server accepts requests immediately and returns cached data while the reasoner warms up.
\end{itemize}

\subsection{Dependencies}

\begin{longtable}{@{}P{3.5cm}P{11cm}@{}}
\toprule
\textbf{Package} & \textbf{Purpose} \\
\midrule
\endhead
\texttt{owlready2} & OWL ontology loading and SPARQL querying via \texttt{default\_world} \\
Java 11+ & Required in \texttt{PATH}; used by owlready2 to invoke the Pellet subprocess \\
\texttt{threading} & \texttt{Lock} objects for \texttt{isSelected} and HC mutation serialisation \\
\texttt{hashlib} & SHA-256 fingerprint of the RDF file for cache invalidation \\
\texttt{json} & Serialisation of the inference snapshot \\
\bottomrule
\end{longtable}

\clearpage
\section{Module-Level Constants}

\subsection{IRI Constants (18)}

All property IRIs used in SPARQL queries are string constants. Two additional IRIs are present compared to \texttt{ontology\_reasoner.py}:

\begin{longtable}{@{}P{5cm}P{9.5cm}@{}}
\toprule
\textbf{Constant} & \textbf{IRI / Role} \\
\midrule
\endhead
\texttt{IRI\_IS\_VERY\_IMP \ldots IRI\_IS\_LESS} & Importance properties inferred by R9--R12 \\
\texttt{IRI\_HAS\_CRIT} & \texttt{hasSpecificCriticality} inferred by R1--R8 \\
\texttt{IRI\_REQUIRES}, \texttt{IRI\_CONCERNS}, \texttt{IRI\_HAS\_SCORE} & Job structure \\
\texttt{IRI\_IS\_EVAL\_JOB}, \texttt{IRI\_IS\_IN\_HC}, \texttt{IRI\_IS\_SELECTED} & Person--job--HC relations \\
\texttt{IRI\_IS\_TRANSL}, \texttt{IRI\_BFQUAL}, \texttt{IRI\_AP1QUAL} & ICF translation and qualifiers \\
\texttt{IRI\_FIRST\_NAME}, \texttt{IRI\_SURNAME} & Person demographics \\
\texttt{IRI\_IS\_DESCRIBED\_BY} & HC $\to$ Descriptor relation \\
\texttt{IRI\_INVOLVES\_ICF} & Descriptor $\to$ ICF code \\
\texttt{IRI\_ICF\_DESCRIPTION} & ICF English description annotation \\
\texttt{IRI\_ONET\_DEFINITION} & O*NET Skill/Ability definition text \\
\bottomrule
\end{longtable}

\subsection{Other Constants}

\begin{longtable}{@{}P{5cm}P{9.5cm}@{}}
\toprule
\textbf{Constant} & \textbf{Value / Role} \\
\midrule
\endhead
\texttt{JS\_RED\_INTERCEPT} & 21.0 --- NOT SUITABLE boundary intercept \\
\texttt{JS\_YELLOW\_INTERCEPT} & 15.5 --- WITH PRECAUTIONS boundary intercept \\
\texttt{JS\_SLOPE} & $-0.5$ --- slope of both boundary lines \\
\texttt{IMPORTANCE\_IRIS} & Dict mapping importance label strings to their IRIs \\
\texttt{SUITABILITY\_COLORS} & Dict mapping suitability labels to hex colour strings \\
\texttt{IRI\_ICF\_CORE\_SET} & IRI of the \texttt{ICF\_Core\_set} class (used for core-set map) \\
\bottomrule
\end{longtable}

\clearpage
\section{Global State --- \texttt{ReasonerState}}

\subsection{Class Definition}

\texttt{ReasonerState} is a \textbf{thread-safe singleton} holding the service readiness status. The module-level instance \texttt{state} is imported by \texttt{main.py} and exposed via the \texttt{/status} endpoint.

\begin{longtable}{@{}P{3cm}P{2.5cm}P{9cm}@{}}
\toprule
\textbf{Attribute} & \textbf{Type} & \textbf{Description} \\
\midrule
\endhead
\texttt{status} & \texttt{str} & One of: \texttt{"loading"}, \texttt{"ready"}, \texttt{"error"} \\
\texttt{message} & \texttt{str} & Human-readable status message \\
\texttt{ontology\_path} & \texttt{str} & Resolved absolute path of the loaded RDF file \\
\texttt{elapsed\_pellet} & \texttt{float|None} & Pellet wall-clock time in seconds \\
\texttt{stats} & \texttt{dict} & Ontology stats: classes, individuals, properties \\
\bottomrule
\end{longtable}

\subsection{Methods}

\begin{longtable}{@{}P{5cm}P{9.5cm}@{}}
\toprule
\textbf{Method} & \textbf{Action} \\
\midrule
\endhead
\texttt{set\_ready(path, elapsed, stats)} & Sets status to \texttt{"ready"} after live Pellet run \\
\texttt{set\_ready\_from\_cache(path, stats)} & Sets status to \texttt{"ready"} after snapshot load \\
\texttt{set\_loading(message)} & Sets status to \texttt{"loading"} \\
\texttt{set\_error(error)} & Sets status to \texttt{"error"} \\
\texttt{is\_ready()} & Returns \texttt{True} if status is \texttt{"ready"} \\
\bottomrule
\end{longtable}

All methods acquire \texttt{self.\_lock} before writing, ensuring thread-safe reads from the FastAPI event loop.

\subsection{Module-Level Variables}

\begin{longtable}{@{}P{5cm}P{9.5cm}@{}}
\toprule
\textbf{Variable} & \textbf{Role} \\
\midrule
\endhead
\texttt{state} & \texttt{ReasonerState} singleton exported to \texttt{main.py} \\
\texttt{\_selection\_lock} & \texttt{threading.Lock} for \texttt{isSelected} mutations \\
\texttt{\_hc\_lock} & \texttt{threading.Lock} for HC mutations \\
\texttt{\_icf\_core\_set\_map} & Dict: ICF code $\to$ sorted list of core-set names \\
\texttt{\_icf\_ind\_cache} & Dict: pure ICF code $\to$ owlready2 individual \\
\texttt{\_snapshot\_cache} & Loaded JSON snapshot dict or \texttt{None} \\
\texttt{\_live\_reasoner\_ready} & \texttt{bool}: True once Pellet has completed \\
\texttt{\_pending\_selected\_worker\_id} & Worker ID to re-select after Pellet completes \\
\texttt{\_loaded\_ontology} & The in-memory owlready2 \texttt{Ontology} object \\
\bottomrule
\end{longtable}

\clearpage
\section{Utility Functions}

\begin{longtable}{@{}P{4cm}P{4cm}P{7cm}@{}}
\toprule
\textbf{Function} & \textbf{Signature} & \textbf{Description} \\
\midrule
\endhead
\texttt{local\_name} & \texttt{(entity) -> str} & Extracts IRI local fragment via \texttt{entity.name} or split on \texttt{\#}/\texttt{/} \\
\texttt{criticality\_label} & \texttt{(cs: int) -> str} & Maps CS value to label (Fig.~2 of paper) \\
\texttt{job\_suitability} & \texttt{(gcs, aisa) -> tuple} & Returns \texttt{(label, hex\_color)} per Fig.~4 \\
\texttt{\_score\_to\_anchor} & \texttt{(score: int) -> int} & Converts O*NET score to anchor $\in \{0,1,2,3\}$ \\
\texttt{\_sparql} & \texttt{(query: str) -> list} & Thin wrapper around \texttt{default\_world.sparql()} \\
\texttt{\_find\_rdf\_file} & \texttt{() -> str} & Auto-detects first \texttt{.rdf/.owl/.ttl/.n3} in service dir or parent \\
\texttt{\_resolve\_rdf\_path} & \texttt{(rdf\_path) -> str} & Resolves path from arg, \texttt{ONTOLOGY\_PATH} env var, or auto-detect \\
\texttt{is\_live\_reasoner\_ready} & \texttt{() -> bool} & Returns \texttt{\_live\_reasoner\_ready} flag \\
\texttt{\_use\_snapshot\_cache} & \texttt{() -> bool} & True when snapshot is loaded and live reasoner is not yet ready \\
\bottomrule
\end{longtable}

\clearpage
\section{Snapshot Cache System}

\subsection{Purpose}

The snapshot cache allows the service to serve \textbf{fully pre-computed read data} on the very first request after startup, before Pellet has finished running (which takes 30--120 seconds). It is a JSON file stored under \texttt{python-service/reasoning\_cache/}.

\subsection{Cache Invalidation}

\subsubsection*{\texttt{\_ontology\_fingerprint(path) -> str}}

Builds a 16-character SHA-256 fingerprint from the absolute path, file size (\texttt{st\_size}), and last modification timestamp (\texttt{st\_mtime\_ns}). The cache filename is \texttt{inference\_\{fingerprint\}.json}, so any change to the RDF file automatically produces a new filename and the old cache is ignored.

\subsection{\texttt{load\_snapshot\_cache(rdf\_path, update\_state) -> bool}}

Loads the JSON snapshot if it exists and is structurally valid. The function checks that the payload contains all 12 required keys:

\begin{longtable}{@{}P{5cm}P{9.5cm}@{}}
\toprule
\textbf{Key} & \textbf{Content} \\
\midrule
\endhead
\texttt{source\_hash} & Fingerprint of the RDF file at save time \\
\texttt{source\_path} & Absolute path of the RDF file \\
\texttt{stats} & Class/individual/property counts \\
\texttt{workers} & All worker dicts (output of \texttt{get\_workers()}) \\
\texttt{jobs} & All job dicts (output of \texttt{get\_jobs()}) \\
\texttt{core\_sets} & All core-set names \\
\texttt{icf\_codes} & All ICF code dicts \\
\texttt{health\_conditions} & Per-worker HC dict \\
\texttt{importance\_by\_job} & Per-job importance summary \\
\texttt{job\_profiles} & Per-job skill profile \\
\texttt{match\_results\_by\_worker} & Per-worker match results \\
\texttt{skill\_details\_by\_worker\_job} & Per-(worker, job) skill detail \\
\bottomrule
\end{longtable}

Returns \texttt{True} and populates \texttt{\_snapshot\_cache} on success. Calls \texttt{state.set\_ready\_from\_cache()} if \texttt{update\_state=True}.

\subsection{\texttt{\_save\_snapshot\_cache(source\_path, elapsed, stats)}}

Iterates over all workers and jobs, temporarily activating each worker with \texttt{set\_selected\_worker} to compute their match results and skill details. Collects all data into the payload dict and writes it atomically (write to \texttt{.tmp}, then \texttt{os.replace}) to avoid corrupt files on crash.

\subsection{\texttt{\_save\_live\_ontology(path)}}

Persists the current in-memory \texttt{\_loaded\_ontology} back to disk via \texttt{onto.save(format="rdfxml")}. Called before re-running Pellet after an HC mutation.

\subsection{\texttt{\_refresh\_reasoning\_artifacts()}}

Called at the end of \texttt{update\_health\_conditions}. Re-runs the full inference pipeline:
\begin{enumerate}[noitemsep]
  \item Save ontology to disk.
  \item Re-run Pellet (\texttt{\_run\_pellet}).
  \item Rebuild ICF core-set map and ICF individual cache.
  \item Update \texttt{\_live\_reasoner\_ready} and \texttt{state}.
  \item Overwrite the JSON snapshot.
  \item Reload snapshot into \texttt{\_snapshot\_cache}.
\end{enumerate}

\clearpage
\section{Phases 1 and 2 --- Load and Reason}

\subsection{\texttt{\_load\_ontology(path) -> Ontology}}

Loads the RDF file into \texttt{owlready2.default\_world} via \texttt{get\_ontology(iri).load()}. Redirects \texttt{stderr} to suppress ICF cycle warnings. Raises \texttt{RuntimeError} on failure.

\subsection{\texttt{\_run\_pellet(onto) -> float|None}}

Runs \texttt{sync\_reasoner\_pellet} with \texttt{infer\_property\_values=True} and \texttt{infer\_data\_property\_values=True}. Applies the same ICF inheritance cycle monkey-patch as \texttt{ontology\_reasoner.py} (patches \texttt{ThingMeta.\_\_setattr\_\_} to swallow \texttt{TypeError: inheritance cycle}, always restoring the original in \texttt{finally}). Parses Pellet's stderr log for elapsed time. Raises \texttt{RuntimeError} on \texttt{UnsupportedClassVersionError} (Java too old) or any other error.

\subsection{\texttt{load\_and\_reason(rdf\_path) -> None}}

The \textbf{public entrypoint} called from \texttt{main.py} in a background thread during the FastAPI lifespan startup. Execution flow:

\begin{enumerate}[noitemsep]
  \item Resolve RDF path (arg $\to$ env var $\to$ auto-detect).
  \item Load ontology; on error call \texttt{state.set\_error()} and return.
  \item Run Pellet; on error call \texttt{state.set\_error()} and return.
  \item Set \texttt{\_live\_reasoner\_ready = True}.
  \item If a worker was pending selection, call \texttt{set\_selected\_worker}.
  \item Call \texttt{state.set\_ready()}.
  \item Build ICF core-set map (\texttt{\_build\_icf\_core\_set\_map}).
  \item Save and reload snapshot cache (best-effort, errors swallowed).
\end{enumerate}

\clearpage
\section{ICF Core-Set Map}

\subsection{\texttt{\_build\_icf\_core\_set\_map()}}

Runs a transitive SPARQL query to find which \texttt{ICF\_Core\_set} subclass each ICF code belongs to:

\begin{lstlisting}[language=SPARQL]
SELECT DISTINCT ?icf ?coreset WHERE {
    ?icf     rdfs:subClassOf+ ?coreset .
    ?coreset rdfs:subClassOf  <ICF_Core_set> .
}
\end{lstlisting}

Populates \texttt{\_icf\_core\_set\_map}: a dict mapping each ICF local name (e.g., \texttt{b110}) to a sorted list of core-set names (e.g., \texttt{["Spinal Cord Injury", "Vocational Rehabilitation"]}).

\subsection{\texttt{get\_all\_core\_sets() -> list[str]}}

Returns a sorted list of all distinct core-set names. Served from \texttt{\_snapshot\_cache["core\_sets"]} when the cache is active.

\clearpage
\section{Query Functions}

\subsection{\texttt{get\_workers() -> list[dict]}}

Returns all \texttt{Person} individuals that have at least one \texttt{isEvaluatedForJob} triple. Each dict:

\begin{longtable}{@{}P{4cm}P{2.5cm}P{8cm}@{}}
\toprule
\textbf{Key} & \textbf{Type} & \textbf{Description} \\
\midrule
\endhead
\texttt{id} & \texttt{str} & Local name of the individual \\
\texttt{first\_name} & \texttt{str} & FOAF first name \\
\texttt{surname} & \texttt{str} & FOAF surname \\
\texttt{is\_selected} & \texttt{bool} & Current value of \texttt{isSelected} \\
\texttt{evaluated\_for\_jobs} & \texttt{list[str]} & Local names of evaluated jobs \\
\bottomrule
\end{longtable}

\subsection{\texttt{get\_jobs() -> list[dict]}}

Returns all Job individuals that have at least one \texttt{requires} triple. Each dict: \texttt{\{id, label\}} where \texttt{label} is the local name with underscores replaced by spaces.

\subsection{\texttt{get\_health\_conditions(worker\_id) -> dict}}

Returns the ICF health condition table for one worker. SPARQL joins: person $\to$ HC $\to$ Descriptor $\to$ ICF code (with \texttt{BFqual} and \texttt{AP1qual}). Handles compound ICF identifiers like \texttt{b1408-AuditoryAttention} by splitting on \texttt{-} and inserting spaces before capitals. Each condition entry:

\begin{longtable}{@{}P{3.5cm}P{2.5cm}P{8.5cm}@{}}
\toprule
\textbf{Key} & \textbf{Type} & \textbf{Description} \\
\midrule
\endhead
\texttt{icf\_code} & \texttt{str} & Pure code, e.g. \texttt{b1408} \\
\texttt{icf\_name} & \texttt{str} & Human-readable name \\
\texttt{description} & \texttt{str} & English description from ontology annotation \\
\texttt{core\_sets} & \texttt{list[str]} & Core-set memberships from \texttt{\_icf\_core\_set\_map} \\
\texttt{bf\_qualifier} & \texttt{int|None} & Body Function qualifier (0--4) \\
\texttt{ap1\_qualifier} & \texttt{int|None} & Activity \& Participation qualifier (0--4) \\
\bottomrule
\end{longtable}

\subsection{\texttt{get\_importance\_summary(job\_id) -> list[dict]}}

Q4 --- reads \texttt{is*ImportantFor} triples (R9--R12) for all or one job. Returns a list of \texttt{\{job\_id, importance\_level, skills: [\{id, score\}]\}} dicts, sorted by job and importance level.

\subsection{\texttt{get\_job\_skill\_profile(job\_id) -> list[dict]}}

Returns every Skill/Ability required by the given job with its raw O*NET score. Pure job-side data (no worker). Used by the React frontend radar chart. Returns \texttt{[\{id, score\}]} sorted by score descending.

\subsection{\texttt{get\_match\_results(worker\_id) -> list[dict]}}

Q1+Q2 --- computes GCS\% and AISA\% for every (Person, Job) pair where \texttt{isSelected=true}. Mirrors \texttt{query\_gcs\_aisa()} from \texttt{ontology\_reasoner.py} but returns JSON-serialisable dicts. Each result:

\begin{longtable}{@{}P{4cm}P{2.5cm}P{8cm}@{}}
\toprule
\textbf{Key} & \textbf{Type} & \textbf{Description} \\
\midrule
\endhead
\texttt{worker\_id} & \texttt{str} & Person local name \\
\texttt{job\_id} & \texttt{str} & Job local name \\
\texttt{gcs\_pct} & \texttt{float} & General Criticality Score \% \\
\texttt{aisa\_pct} & \texttt{float} & Amount Impaired SkAb \% \\
\texttt{n\_total} & \texttt{int} & Total SkAb count \\
\texttt{n\_critical} & \texttt{int} & SkAb with CS $>$ 0 \\
\texttt{suitability} & \texttt{str} & Suitability label \\
\texttt{suitability\_color} & \texttt{str} & Hex colour string \\
\bottomrule
\end{longtable}

Falls back to \texttt{\_match\_results\_fallback} (uses \texttt{hasSpecificCriticality}) if the primary SPARQL join returns no rows.

\subsection{\texttt{get\_skill\_detail(worker\_id, job\_id) -> dict}}

Q3 --- full Skill/Ability breakdown for one (worker, job) pair. Queries \texttt{hasSpecificCriticality} joined with \texttt{hasScore} and the O*NET definition. Returns \texttt{\{worker\_id, job\_id, skills: [...]}\} with skills sorted by CS descending.

Each skill entry:

\begin{longtable}{@{}P{4cm}P{2.5cm}P{8cm}@{}}
\toprule
\textbf{Key} & \textbf{Type} & \textbf{Description} \\
\midrule
\endhead
\texttt{id} & \texttt{str} & SkAb local name \\
\texttt{score} & \texttt{int} & O*NET importance score \\
\texttt{importance\_label} & \texttt{str} & e.g. \texttt{isVeryImportantFor} \\
\texttt{anchor} & \texttt{int} & 0--3 \\
\texttt{qualifier} & \texttt{int} & ICF qualifier derived from $\text{CS} \div \text{anchor}$ \\
\texttt{cs} & \texttt{int} & Criticality Score \\
\texttt{cs\_normalized} & \texttt{float} & CS / 12 \\
\texttt{criticality\_label} & \texttt{str} & e.g. EXTREMELY CRITICAL \\
\texttt{description} & \texttt{str} & O*NET definition text \\
\bottomrule
\end{longtable}

\subsection{\texttt{get\_all\_icf\_codes() -> list[dict]}}

Returns every ICF code that appears in any \texttt{involvesICFCode} triple. Used to populate the selection table in the Modify-Health-Condition wizard. Merges multiple ontology individuals that collapse to the same pure code. Each entry includes \texttt{icf\_code}, \texttt{icf\_name}, \texttt{description}, \texttt{category} (Body Functions / Activities and Participation / Body Structures / Environmental Factors), \texttt{core\_sets}, and \texttt{iri}.

\clearpage
\section{Mutation Functions}

\subsection{\texttt{set\_selected\_worker(worker\_id) -> dict}}

Atomically deselects the currently selected worker and selects the given one. Uses \texttt{\_selection\_lock} to prevent concurrent corruption.

\subsubsection*{Live mode (Pellet ready)}
Scans all individuals for those with \texttt{isSelected=True}, sets them to \texttt{False}, then sets the target to \texttt{True}. Returns \texttt{\{"previous": old\_id, "selected": worker\_id\}}.

\subsubsection*{Cache mode (Pellet still warming)}
Updates the in-memory \texttt{\_snapshot\_cache["workers"]} list directly and stores \texttt{worker\_id} in \texttt{\_pending\_selected\_worker\_id}. When Pellet completes, \texttt{load\_and\_reason} replays the selection on the live ontology.

\subsection{\texttt{update\_health\_conditions(worker\_id, changes) -> dict}}

Applies a batch of HC changes. Requires \texttt{\_live\_reasoner\_ready = True} (raises \texttt{RuntimeError} otherwise). Protected by \texttt{\_hc\_lock}.

\subsubsection*{Change dict schema}
\begin{lstlisting}
{
    "icf_code" : str,           # e.g. "b1408"
    "action"   : "add" | "remove" | "modify",
    "qualifier": int | None     # 0-4; required for add/modify
}
\end{lstlisting}

\subsubsection*{ADD action}
\begin{enumerate}[noitemsep]
  \item Resolve ICF individual via \texttt{\_resolve\_icf\_individual}.
  \item If code already in \texttt{desc\_map}: treat as modify (update qualifier only).
  \item Otherwise: call \texttt{\_get\_or\_create\_hc()} to get/create the HC individual.
  \item Instantiate a new \texttt{HC\_Descriptor} individual.
  \item Set \texttt{involvesICFCode}, \texttt{BFqual} (b-codes) or \texttt{AP1qual} (d-codes), clear the other qualifier.
  \item Append descriptor to HC via \texttt{isDescribedBy}.
\end{enumerate}

\subsubsection*{REMOVE action}
\begin{enumerate}[noitemsep]
  \item Locate the Descriptor in \texttt{desc\_map}.
  \item Remove it from \texttt{HC.isDescribedBy}.
  \item Call \texttt{owlready2.destroy\_entity(des\_ind)} to remove all triples.
  \item If HC has no remaining descriptors: destroy the HC individual and unlink from person.
\end{enumerate}

\subsubsection*{MODIFY action}
\begin{enumerate}[noitemsep]
  \item Locate the Descriptor.
  \item Set the appropriate qualifier property (\texttt{BFqual} or \texttt{AP1qual}).
  \item Clear the other qualifier to prevent stale data.
\end{enumerate}

\subsubsection*{Post-mutation}
After all changes, calls \texttt{\_refresh\_reasoning\_artifacts()}: saves the RDF to disk, re-runs Pellet, rebuilds caches, and overwrites the JSON snapshot.

\subsubsection*{Return value}
\begin{lstlisting}
{ "worker_id": str, "added": int, "removed": int,
  "modified": int, "errors": [str, ...] }
\end{lstlisting}

\subsection{ICF Individual Resolution --- \texttt{\_resolve\_icf\_individual(icf\_code)}}

Finds the owlready2 individual for a pure ICF code using a four-step fallback:
\begin{enumerate}[noitemsep]
  \item Check \texttt{\_icf\_ind\_cache} (O(1) dict lookup).
  \item \texttt{default\_world.search\_one(iri="*\#" + icf\_code + "-*")} (wildcard prefix).
  \item \texttt{default\_world.search\_one(iri="*\#" + icf\_code)} (exact match).
  \item Linear scan over all individuals (last resort).
\end{enumerate}

\clearpage
\section{Internal Helper Functions}

\begin{longtable}{@{}P{5.5cm}P{9cm}@{}}
\toprule
\textbf{Function} & \textbf{Role} \\
\midrule
\endhead
\texttt{\_build\_icf\_ind\_cache()} & Scans \texttt{involvesICFCode} triples; populates \texttt{\_icf\_ind\_cache} \\
\texttt{\_is\_ap1\_code(icf\_code)} & Returns \texttt{True} for d-prefix codes (AP1qual) \\
\texttt{\_match\_results\_fallback(worker\_id)} & Computes GCS\%/AISA\% from \texttt{hasSpecificCriticality} when primary query fails \\
\texttt{\_compute\_metrics(agg)} & Converts aggregated \texttt{\{skab: \{score, qual\}\}} to GCS\%/AISA\% result list \\
\texttt{\_get\_or\_create\_hc()} & Inner closure in \texttt{update\_health\_conditions}: gets the worker's HC or creates a new one \\
\texttt{\_snapshot\_cache\_path(path)} & Returns the absolute path of the cache JSON file \\
\bottomrule
\end{longtable}

\clearpage
\section{Function Reference Summary}

\begin{longtable}{@{}P{6.5cm}P{3cm}P{5.5cm}@{}}
\toprule
\textbf{Function} & \textbf{Returns} & \textbf{Public?} \\
\midrule
\endhead
\texttt{load\_and\_reason(rdf\_path)} & \texttt{None} & Yes (called by \texttt{main.py}) \\
\texttt{load\_snapshot\_cache(path, update)} & \texttt{bool} & Yes \\
\texttt{is\_live\_reasoner\_ready()} & \texttt{bool} & Yes \\
\texttt{get\_workers()} & \texttt{list[dict]} & Yes \\
\texttt{get\_jobs()} & \texttt{list[dict]} & Yes \\
\texttt{get\_health\_conditions(worker\_id)} & \texttt{dict} & Yes \\
\texttt{get\_importance\_summary(job\_id)} & \texttt{list[dict]} & Yes \\
\texttt{get\_job\_skill\_profile(job\_id)} & \texttt{list[dict]} & Yes \\
\texttt{get\_match\_results(worker\_id)} & \texttt{list[dict]} & Yes \\
\texttt{get\_skill\_detail(worker, job)} & \texttt{dict} & Yes \\
\texttt{get\_all\_icf\_codes()} & \texttt{list[dict]} & Yes \\
\texttt{get\_all\_core\_sets()} & \texttt{list[str]} & Yes \\
\texttt{set\_selected\_worker(worker\_id)} & \texttt{dict} & Yes \\
\texttt{update\_health\_conditions(w, changes)} & \texttt{dict} & Yes \\
\texttt{local\_name(entity)} & \texttt{str} & Yes \\
\texttt{criticality\_label(cs)} & \texttt{str} & Yes \\
\texttt{job\_suitability(gcs, aisa)} & \texttt{tuple} & Yes \\
\texttt{\_load\_ontology(path)} & Ontology & Internal \\
\texttt{\_run\_pellet(onto)} & \texttt{float|None} & Internal \\
\texttt{\_build\_icf\_core\_set\_map()} & \texttt{None} & Internal \\
\texttt{\_build\_icf\_ind\_cache()} & \texttt{None} & Internal \\
\texttt{\_resolve\_icf\_individual(code)} & object|None & Internal \\
\texttt{\_refresh\_reasoning\_artifacts()} & \texttt{None} & Internal \\
\texttt{\_save\_snapshot\_cache(...)} & \texttt{None} & Internal \\
\texttt{\_save\_live\_ontology(path)} & \texttt{None} & Internal \\
\texttt{\_match\_results\_fallback(worker\_id)} & \texttt{list[dict]} & Internal \\
\texttt{\_compute\_metrics(agg)} & \texttt{list[dict]} & Internal \\
\bottomrule
\end{longtable}

\clearpage
\section{Known Issues and Design Notes}

\begin{tcolorbox}[colback=noteblue!40, colframe=blue!60!black,
                  title=Note: Two-tier serving model]
When the service starts, \texttt{load\_snapshot\_cache()} is called first. If a valid cache exists, \texttt{state} is immediately set to \texttt{"ready"} and all read endpoints can respond. Simultaneously, \texttt{load\_and\_reason()} runs in a background thread. Once Pellet completes, \texttt{\_live\_reasoner\_ready} flips to \texttt{True} and \texttt{\_use\_snapshot\_cache()} returns \texttt{False}: all subsequent reads go directly to \texttt{default\_world} via SPARQL.
\end{tcolorbox}

\vspace{0.6em}

\begin{tcolorbox}[colback=warnred!40, colframe=red!60!black,
                  title=Warning: HC mutations require the live reasoner]
\texttt{update\_health\_conditions} raises \texttt{RuntimeError} if \texttt{\_live\_reasoner\_ready} is \texttt{False}. Mutations cannot be applied to the snapshot --- they require the live in-memory owlready2 graph. The frontend should poll \texttt{/status} until \texttt{status="ready"} and \texttt{elapsed\_pellet} is non-null before enabling the Modify-Health-Condition wizard.
\end{tcolorbox}

\vspace{0.6em}

\begin{tcolorbox}[colback=importantbg!60, colframe=green!50!black,
                  title=Important: Atomic cache write]
The snapshot JSON is written to a \texttt{.tmp} file first, then renamed via \texttt{os.replace()}. This is an atomic operation on all supported platforms and prevents serving a partially written cache if the process is killed during a write.
\end{tcolorbox}

\end{document}
